\documentclass{beamer}
\usetheme{Madrid}
\usecolortheme{dolphin}
\definecolor{charcoal}{RGB}{54,69,79}
\setbeamercolor{structure}{fg=charcoal}
\setbeamercolor{palette primary}{bg=charcoal,fg=white}
\setbeamercolor{palette secondary}{bg=charcoal!90,fg=white}
\setbeamercolor{palette tertiary}{bg=charcoal!80,fg=white}
\setbeamercolor{palette quaternary}{bg=charcoal!70,fg=white}
\setbeamercolor{frametitle}{bg=charcoal,fg=white}
\setbeamercolor{title}{bg=charcoal,fg=white}
\usepackage{hyperref}
\usepackage{graphicx}
\usepackage{booktabs}
\usepackage{listings}
\usepackage{color}
\definecolor{dkgreen}{rgb}{0,0.6,0}
\definecolor{gray}{rgb}{0.5,0.5,0.5}
\definecolor{mauve}{rgb}{0.58,0,0.82}
\lstset{
  language=Java,
  aboveskip=3mm,
  belowskip=3mm,
  showstringspaces=false,
  columns=flexible,
  basicstyle={\small\ttfamily},
  numbers=none,
  keywordstyle=\color{blue},
  commentstyle=\color{dkgreen},
  stringstyle=\color{mauve},
  breaklines=true,
  breakatwhitespace=true,
  tabsize=2
}
\title[Multi-Dim Arrays]{Multi-Dimensional Arrays}
\subtitle{COP2250: Java Programming}
\author{Kevin Pyatt, Ph.D.}
\institute{State College of Florida \\ Pyatt Labs}
\date{Week 9}
\begin{document}
\begin{frame}
  \titlepage
\end{frame}
\begin{frame}{Today's Objectives}
\begin{itemize}
  \item Declare and create two-dimensional arrays
  \item Access elements using row and column indices
  \item Use \texttt{m.length} and \texttt{m[0].length} correctly
  \item Traverse a matrix with nested loops
  \item Sum rows and columns using fixed-index traversal
  \item Pass 2D arrays to methods
\end{itemize}
\end{frame}
\section{From 1D to 2D}
\begin{frame}[fragile]{From One Row to a Grid}
\textbf{Last week:} one row of values
\begin{lstlisting}
double[] scores = new double[5];
\end{lstlisting}
\vspace{0.3cm}
\textbf{This week:} a grid of values (matrix)
\begin{lstlisting}
double[][] m = new double[3][4];
\end{lstlisting}
\vspace{0.3cm}
\begin{tabular}{ccccc}
 & col 0 & col 1 & col 2 & col 3 \\
\textbf{row 0} & 1.0 & 2.0 & 3.0 & 4.0 \\
\textbf{row 1} & 5.0 & 6.0 & 7.0 & 8.0 \\
\textbf{row 2} & 9.0 & 10.0 & 11.0 & 12.0 \\
\end{tabular}
\vspace{0.3cm}
\textit{Think spreadsheet. Rows and columns. Row first, column second. Always.}
\end{frame}
\section{Array Basics}
\begin{frame}[fragile]{Declaring and Creating 2D Arrays}
\begin{lstlisting}
// Declare and create (all values default to 0.0)
double[][] m = new double[3][4];

// Declare and initialize with values
int[][] grid = {
    {1, 2, 3},
    {4, 5, 6},
    {7, 8, 9}
};
\end{lstlisting}
\vspace{0.3cm}
\begin{tabular}{ll}
\texttt{m.length} & Number of \textbf{rows} (3) \\
\texttt{m[0].length} & Number of \textbf{columns} (4) \\
\texttt{m[i][j]} & Element at row \texttt{i}, column \texttt{j} \\
\end{tabular}
\vspace{0.3cm}
\textit{Row first. Column second. Every time.}
\end{frame}
\section{Processing Arrays}
\begin{frame}[fragile]{Nested Loops for 2D Arrays}
\textbf{One loop for one dimension. Two loops for two.}
\begin{lstlisting}
for (int i = 0; i < m.length; i++) {        // rows
    for (int j = 0; j < m[0].length; j++) { // columns
        System.out.printf("%6.1f", m[i][j]);
    }
    System.out.println();
}
\end{lstlisting}
\vspace{0.3cm}
\begin{itemize}
  \item Outer loop = rows (\texttt{i})
  \item Inner loop = columns (\texttt{j})
  \item Element = \texttt{m[i][j]}
\end{itemize}
\vspace{0.2cm}
\textit{Every time you see a 2D array, think: nested loop.}
\end{frame}
\begin{frame}[fragile]{Traversing a Row vs a Column}
\textbf{Sum a row} --- fix the row, loop the columns:
\begin{lstlisting}
double sum = 0;
for (int j = 0; j < m[0].length; j++) {
    sum += m[rowIndex][j];
}
\end{lstlisting}
\vspace{0.3cm}
\textbf{Sum a column} --- fix the column, loop the rows:
\begin{lstlisting}
double sum = 0;
for (int i = 0; i < m.length; i++) {
    sum += m[i][columnIndex];
}
\end{lstlisting}
\vspace{0.3cm}
\textbf{Key rule:} To sum a column, the column index is \textbf{FIXED}. The row index \textbf{loops}.\\
\vspace{0.2cm}
\textit{Write it down. This is your assignment.}
\end{frame}
\begin{frame}[fragile]{The sumColumn Method}
\textbf{Required header (use exactly):}
\begin{lstlisting}
public static double sumColumn(double[][] m, int columnIndex)
\end{lstlisting}
\vspace{0.3cm}
\textbf{Implementation:}
\begin{lstlisting}
public static double sumColumn(double[][] m, int columnIndex) {
    double sum = 0;
    for (int i = 0; i < m.length; i++) {
        sum += m[i][columnIndex];
    }
    return sum;
}
\end{lstlisting}
\vspace{0.3cm}
\textbf{Usage in main:}
\begin{lstlisting}
for (int col = 0; col < m[0].length; col++) {
    System.out.println("Sum of column " + col
        + " is " + sumColumn(m, col));
}
\end{lstlisting}
\end{frame}
\begin{frame}[fragile]{Common Mistakes}
\textbf{1. Using m.length for columns:}
\begin{lstlisting}
for (int j = 0; j < m.length; j++)      // WRONG
for (int j = 0; j < m[0].length; j++)   // CORRECT
\end{lstlisting}
\vspace{0.2cm}
\textbf{2. Swapping row and column in access:}
\begin{lstlisting}
sum += m[columnIndex][i];  // WRONG
sum += m[i][columnIndex];  // CORRECT
\end{lstlisting}
\vspace{0.2cm}
\textbf{3. Initializing max to 0:}
\begin{lstlisting}
double max = 0;       // WRONG if all values negative
double max = m[0][0]; // CORRECT
\end{lstlisting}
\vspace{0.2cm}
\textbf{4. Off-by-one in nested loop:}
\begin{lstlisting}
for (int i = 0; i <= m.length; i++)  // CRASH
for (int i = 0; i < m.length; i++)   // CORRECT
\end{lstlisting}
\end{frame}
\begin{frame}{Assignment 8: Sum of Columns}
\textbf{Exercise 8.1:} Write a program that:
\begin{itemize}
  \item Implements \texttt{sumColumn(double[][] m, int columnIndex)}
  \item Reads a 3-by-4 matrix from the user
  \item Displays the sum of each column
\end{itemize}
\vspace{0.4cm}
\textbf{Sample output:}\\
\texttt{Sum of column 0 is 15.0}\\
\texttt{Sum of column 1 is 18.0}\\
\texttt{Sum of column 2 is 21.0}\\
\texttt{Sum of column 3 is 24.0}
\vspace{0.4cm}
\textbf{Lab: MatrixPractice}
\begin{itemize}
  \item Build matrix utility methods: fill, print, sumRow, sumColumn, total, max
\end{itemize}
\vspace{0.3cm}
\textbf{Next Week:} Objects and Classes (Chapter 9)
\end{frame}
\end{document}
